% SYN-ACK WIRE -- user manual for the Mattel Intellivision, with a FujiNet cartridge
% Build: make   (or make in docs/: the rulebook and every manual)
%
% The style, the theme macros and the furniture the shared rules text needs are in ../../style/synack.sty;
% the rules text is ../../rules/*.tex; the Intellivision's names for things (its keys) are in platform.tex.
% Screens are in screens/: dumps of the Intellivision client's screen, drawn by client/tools/render_intv.py
% (client/tools/capture.sh with PLATFORM=intv makes them, from the autoplay ROM in jzIntv).
\documentclass[12pt,letterpaper]{article}

\usepackage{synack}
% The Intellivision's names for things, for the shared style and text (../../style/synack.sty, ../../rules/).
\renewcommand{\DocKind}{USER MANUAL // INTELLIVISION}
\renewcommand{\PlatformShort}{Intellivision}               % as in "on the Intellivision"
\newcommand{\PlatformLong}{Mattel Intellivision}           % the title page
\newcommand{\PlatformMachines}{INTELLIVISION \quad // \quad FUJINET INTELLIVISION CARTRIDGE}   % the title page, the screen chapter
\newcommand{\SignalLost}{the top line turns red and reads \code{NO SIGNAL}}  % the link is quiet
\renewcommand{\KeyReturn}{ENTER}                           % the keypad's key that confirms
\renewcommand{\KeyEsc}{LOWER LEFT, 3}                      % the options button, then the options' 3
\renewcommand{\typeword}{dial}                             % a handle goes in with a letter dial (or the ECS)
\renewcommand{\Typeword}{Dial}
\renewcommand{\HostAsk}{SEATS? \Rule{minplayers}-\Rule{maxplayers}}
\renewcommand{\LogOffAsk}{\MakeUppercase{\V{logoff}}? 1:YES}
\renewcommand{\KeyYes}{1}
\renewcommand{\TooLate}{TOO LATE: SEAT TAKEN}

% Where the screen sits in a capture, in image pixels (see \capgeom): client/tools/render_intv.py draws the
% 20 x 12 cards inside an 8-pixel border, each card 32 x 32 image pixels, so a callout's (column, row) is the card's.
\makeatletter
\newcommand{\rendered}{\capgeom{32}{16}{32}{32}{704}{416}}
\makeatother

\graphicspath{{screens/}}

\makeatletter
% The rules' captures are staged games (server/stage.py, scenarios.toml) on the board of the drawing
% under "1. Spike a node", so the drawing, the text and the screen tell one story.
\expandafter\def\csname shot@spike\endcsname{\screen{spike.png}{CAPTURE \quad \MakeUppercase{\V{spike}} A \MakeUppercase{\V{node}}}%
  {The drawing above, as ZER0COOL sees it. The yellow digits \code{1}--\code{4} sit on its four targets: they are
  the \Vp{exploit} in your \V{kit} that run there, as the line under the \V{grid} lists them. Press the
  \V{exploit}'s number on the keypad to run it.}}
\expandafter\def\csname shot@broker\endcsname{\screen{infiltration.png}{CAPTURE \quad \MakeUppercase{\V{infiltration}}}%
  {Target 2: ZER0COOL ran \node{C6}, next to the lone \V{spike} on \node{B6}. In the panel beside the \V{grid}, a
  yellow digit stands by each syndicate with no \V{hive}: press one's digit on the keypad to \V{broker} the new
  \V{hive} to it.}}
\expandafter\def\csname shot@hivejack\endcsname{\screen{tied-root.png}{CAPTURE \quad \MakeUppercase{\V{hivejack}}: A TIED \MakeUppercase{\V{root}}}%
  {The example below, as ZER0COOL sees it. The spike on \node{G3} joined \Syn{5}{name} (4 \Vp{node}) and \Syn{3}{name}
  (3). The line above the question says which \V{hive} took which and how many of the \V{absorbed} \Vp{coin} you
  hold: key in how many to \V{cashout} and press \key{ENTER}, then how many to \V{swap}, an even number, and
  \key{ENTER} again.}}
\expandafter\def\csname shot@buy\endcsname{\screen{spread-buy.png}{CAPTURE \quad BUY \MakeUppercase{\V{coin}}}%
  {Target 3: \node{E1} spread \Syn{5}{name} to 5 \Vp{node}, and its price is already up. While you buy, the
  panel beside the \V{grid} shows each syndicate's price, a keypad digit by each with a \V{hive}. Each digit buys one
  of its \V{coin}, up to \Rule{maxbuy} in any mix; \key{CLEAR} takes one back, \key{ENTER} buys.}}
\expandafter\def\csname shot@nullroute\endcsname{\screen{nullroute.png}{CAPTURE \quad \MakeUppercase{\V{lockdown}}: THE QUESTION}%
  {ZER0COOL's \V{spike} on \node{H2} spread \Syn{2}{name} to \Rule{safesize} \Vp{node}: \V{safe}, and with
  \Syn{5}{name} and \Syn{0}{name} already \V{safe}, every \V{hive} is. With too little \Game{currency} for any
  \V{coin}, the prompt is the question alone: \key{1} calls the \V{nullroute}, \key{0} or \key{ENTER} plays on.}}
\expandafter\def\csname shot@lobby\endcsname{\screen{lobby-join.png}{CAPTURE \quad JOINING A LOBBY}%
  {ZER0COOL hosts and MAMBA is waiting; the Intellivision has just joined, and shows the seats. The host clock
  started when MAMBA arrived, and every newcomer is told how long the host has left.}}
\expandafter\def\csname shot@final\endcsname{\screen{game-over.png}{CAPTURE \quad THE END OF A GAME}%
  {ZER0COOL called it. \VC{root} and \V{sudo} were paid for every \V{hive}, all \V{coin} was sold, and the top
  \V{player} shows above the question. The table is a lobby again: the host's prompt asks for the seats of a
  rematch.}}
\makeatother


\begin{document}

% ================================================================== title
\begin{titlepage}
\machine
\vspace*{1.6in}
{\color{amber}\rule{\linewidth}{2pt}}\par\vspace{14pt}
{\fontsize{44}{48}\selectfont \game\par}\vspace{10pt}
{\fontsize{12}{16}\selectfont\color{steel} SEVEN SYNDICATES. ONE REGION. EVERY \MakeUppercase{\V{node}} FOR SALE.\par}
\vspace{14pt}{\color{amber}\rule{\linewidth}{2pt}}\par
\vspace{0.35in}
{\fontsize{24}{30}\selectfont USER MANUAL\par}\vspace{4pt}
{\fontsize{14}{18}\selectfont\color{amber} FOR THE \MakeUppercase{\PlatformLong}\par}
\vspace{0.4in}
{\fontsize{12}{16}\selectfont \PlatformMachines\par
\Rule{minplayers}-\Rule{maxplayers} \MakeUppercase{\Vp{player}} \quad // \quad BOTS FILL EMPTY SEATS\par}
\vfill
{\fontsize{8.5}{11}\selectfont\color{steel}
AN UNOFFICIAL CYBERPUNK RE-SKIN OF ACQUIRE BY SID SACKSON.\par
ACQUIRE IS A TRADEMARK OF HASBRO, INC. NOT AFFILIATED WITH OR ENDORSED BY HASBRO.\par}
\end{titlepage}

% ================================================================== 01
\renewcommand{\secstatus}{THE SETTING \quad // \quad THE OBJECT \quad // \quad THE END}
\chapter{The Network Grid}
\begin{multicols}{2}
\rules{opening}
\end{multicols}

% ================================================================== 02
\renewcommand{\secstatus}{REGION \quad // \quad STATUS \quad // \quad HANDLE}
\chapter{\VC{jackin}}
\game\ runs on an Intellivision with a FujiNet Intellivision cartridge. Pick a \game\ server in the FujiNet lobby,
which loads the game itself, or in FujiNet's CONFIG mount \code{syn-ack-wire-intv.rom} from the host
\code{syn-ack-wire.fly.dev} and boot it. The lobby passes the server you pick through FujiNet's SD card: without
one, CHIBA-GRID-7 is the server you join.

The title dials the server as soon as it appears. Dial in a handle, up to 8 characters: the disc up and down
changes the letter the dial is on, right and left move along the handle, \key{CLEAR} deletes a letter. Then press
\key{ENTER}. From the FujiNet lobby, the handle starts as your lobby name. With an ECS keyboard attached, type it.

\begin{capture}{title.png}{CAPTURE 01 \quad TITLE SCREEN \quad LOBBY OPEN}
  \callout{1}{L}{4}{3}
  \callout{2}{R}{14}{7}
  \callout{3}{L}{1}{8}
  \callout{4}{R}{7}{10}
  \callout{5}{L}{1}{11}
\end{capture}

\begin{legend}
\li{1}{\textbf{The logo.} Its letters are static, as on a TV with no station.}
\li{2}{\textbf{Region.} The server's \V{grid}: one game at a time, up to six seats.}
\li{3}{\textbf{Status.} Whether you can join, and how: \code{LOBBY OPEN}, \code{IN PROGRESS}, \code{TABLE FULL}.
It stays up to date while you dial.}
\li{4}{\textbf{Handle.} Your name at the table. The letter the dial is on shows yellow; past the end, a yellow line
takes a new letter.}
\li{5}{\textbf{The keys,} one line at a time: \key{ENTER} \VC{jackin}; the top buttons open the instructions; the
lower left button the options (move hints, sound); the disc dials. A refusal, such as a handle already taken,
shows here in red.}
\end{legend}

% ================================================================== 03
\renewcommand{\secstatus}{20 COLUMNS \quad // \quad THREE VIEWS \quad // \quad COLOUR}
\chapter{The Screen}
\begin{sloppypar}
The Intellivision shows 20 columns of 12 lines, half the other machines' each way, so \game\ splits its screen
into three views. The disc (left or right), or the lower right button, steps through them: the \V{grid}, the
seats, and the \V{coin}. The bottom two lines are the same in all three, and turn purple while the decision is
yours; a new question brings the \V{grid} back. Colour does the work inverse does elsewhere: each syndicate's
\Vp{hive} in its own colour, yours in yellow. This is early in a game, and it's your turn: pick an \V{exploit} to
run.
\end{sloppypar}

\begin{capture}{spike-early.png}{CAPTURE 02 \quad \MakeUppercase{\V{spike}} PROMPT \quad 3 SEATS \quad EARLY GAME}
  \callout{1}{L}{2}{0}
  \callout{2}{L}{1}{6}
  \callout{3}{L}{8}{3}
  \callout{4}{R}{14}{6}
  \callout{5}{R}{19}{6}
  \callout{6}{R}{17}{8}
  \callout{7}{R}{14}{9}
  \callout{8}{L}{1}{10}
  \callout{9}{R}{18}{11}
\end{capture}

\begin{legend}
\li{1}{\textbf{The \V{grid}'s columns}, and the panel's heading. It turns red and reads \code{NO SIGNAL} when the
server goes quiet.}
\li{2}{\textbf{The \V{grid}.} A letter on a colour is a \V{hive}, in its syndicate's colour; \node{\Glyph{spike}}
is a lone \V{spike}. The last move blinks for a few seconds.}
\li{3}{\textbf{Move hints.} A yellow digit marks a \V{node} the \V{exploit} with that number can run on; a red
one, an \V{exploit} that can't run there yet.}
\li{4}{\textbf{\VCp{hive}.} Each syndicate's letter and \V{hive} size (\V{safe}: black on white). While a question
asks for a syndicate, its keypad digit stands beside it; while you buy, its price.}
\li{5}{\textbf{Your \V{coin}}, in yellow, by each syndicate.}
\li{6}{\textbf{Your \Game{currency}.}}
\li{7}{\textbf{Last move.} The \V{node} \V{spike}d most recently.}
\li{8}{\textbf{The info line.} Here your \V{kit}: each \V{exploit}'s number and the \V{node} it runs on, red if it
can't run now. Otherwise what the question is about, a refusal, or the newest log line, which scrolls when it is
long.}
\li{9}{\textbf{The prompt.} What you're asked, with seconds left on the shot clock.}
\end{legend}

\subsection{The other views}
\begin{description}[leftmargin=0pt,labelsep=0.5em]
  \item[The seats.] The region on the top line; each seat's handle and \Game{currency}: \code{>} marks you,
    \code{*} a bot, \code{?} someone away, and the handle whose turn it is shows green. Under them, the log's last
    three lines; a long one goes on, indented, on the next.
  \item[The \V{coin}.] Who holds how many of each syndicate's \V{coin}, handles cut to 5 letters, yours in yellow;
    then how many the \V{bank} has left, each \V{hive}'s size, and while you buy, the prices in hundreds.
\end{description}

% ================================================================== 04
\renewcommand{\secstatus}{THE KEYPAD \quad // \quad THE DISC AND BUTTONS \quad // \quad THE ECS}
\chapter{The Controller}
\game\ is played with the left hand controller: the keypad answers, the disc and the side buttons get around.
Each question on the prompt line names the keys it takes. The game checks an answer before sending it and says on
the info line what's wrong, such as \code{NOT ENOUGH NEW YEN}. An ECS keyboard, if one is attached, works too.

\subsection{Any time}
\begin{keytable}
\kr{\key{TOP}}{Open the instructions, on the title screen or in the middle of a game: either top button.}
\kr{\key{LOWER LEFT}}{The options, over the bottom two lines: \key{1} move hints on or off (the digits on the \V{grid} that show where each \V{exploit} in your \V{kit} can run), \key{2} sound on or off (a chord tells you it's on), and in a game \key{3} \V{logoff}. Any other key closes them.}
\kr{\key{LOWER RIGHT}}{The next view: the \V{grid}, the seats, the \V{coin}. The disc, left or right, steps through them too.}
\end{keytable}

\subsection{The title screen}
\begin{keytable}
\kr{The disc}{Up and down change the letter the dial is on: A--Z, 0--9, then the signs a handle can have. Right and left move along the handle; past its end, up or down adds a letter.}
\kr{\key{CLEAR}}{Delete the letter the dial is on.}
\kr{\key{0}--\key{9}}{Add a digit to the handle.}
\kr{\key{ENTER}}{\VC{jackin}.}
\end{keytable}

\subsection{The instructions}
\begin{keytable}
\kr{The disc}{Right: the next page. Left: the previous one.}
\kr{A button}{Back to where you were. \key{ENTER} and \key{CLEAR} do the same.}
\end{keytable}

\subsection{Hosting}
\begin{keytable}
\kr{\key{\Rule{minplayers}}--\key{\Rule{maxplayers}}}{Start the game with that many seats. Bots take the seats nobody has joined.}
\end{keytable}

\subsection{Your turn}
\begin{keytable}
\kr{\key{1}--\key{6}}{\VC{spike}: run that \V{exploit} in your \V{kit}. A red one can't run now; the info line says why.}
\kr{\key{1}--\key{7}}{\VC{infiltration}: the syndicate to \V{broker} the new \V{hive} to, by its digit in the panel (\Syn{0}{letter} is 1, \Syn{6}{letter} is 7). A tied \V{hivejack}: the \V{hive} that survives, or is \V{absorbed} next.}
\kr{\key{1}--\key{7}}{Buy \V{coin}: each digit one \V{coin} of that syndicate, up to \Rule{maxbuy} in any mix.}
\kr{\key{CLEAR}}{Take back the last \V{coin} in a buy.}
\kr{\key{ENTER}}{Confirm the buy. \key{ENTER} alone buys nothing.}
\kr{\key{1}\ \ \key{0}}{When the game may end: \key{1} calls the \V{nullroute} now, \key{0} (or \key{ENTER}) plays on.}
\end{keytable}

\subsection{When your \V{hive} is \V{absorbed}}
\begin{keytable}
\kr{Digits, \key{ENTER}}{First, how many \Vp{coin} to \V{cashout}; then how many to \V{swap}, an even number.
The rest you \V{hodl}. \key{CLEAR} takes back a digit.}
\end{keytable}

\subsection{Leaving}
\begin{keytable}
\kr{\key{LOWER LEFT}\ \ \key{3}}{\VC{logoff} for good, after \key{1} for yes. A bot takes over your seat at once, and you're back on the title screen. Any other key stays.}
\end{keytable}

\subsection{The ECS keyboard}
\begin{keytable}
\kr{Letters, digits}{Type the handle; at a question, a syndicate's letter or a number, as the keypad's digits would.}
\kr{\key{RTN}}{As \key{ENTER}.}
\kr{\key{$\leftarrow$}}{As \key{CLEAR}: delete, take back.}
\kr{\key{Y}\ \ \key{N}}{Yes and no.}
\kr{\key{?}}{The instructions.}
\kr{\key{ESC}}{The options.}
\kr{\key{$\rightarrow$}\ \ \key{$\uparrow$}\ \ \key{$\downarrow$}}{The next view, or the one before.}
\end{keytable}

% ================================================================== 05
\renewcommand{\secstatus}{THE \MakeUppercase{\V{grid}} \quad // \quad THE SYNDICATES \quad // \quad SETUP}
\chapter{What's in the Game}
\begin{multicols}{2}
\rules{contents}
\end{multicols}

\subsection{The Seven Syndicates}
\syntable

\subsection{Setup}
\begin{multicols}{2}
\rules{setup}
\end{multicols}

% ================================================================== 06
\renewcommand{\secstatus}{\MakeUppercase{\V{spike}} \quad // \quad BUY \quad // \quad DRAW}
\chapter{Playing a Turn}
\begin{multicols}{2}
\rules{turn}
\end{multicols}

% ================================================================== 07
\renewcommand{\secstatus}{\MakeUppercase{\V{totalcompromise}} \quad // \quad \MakeUppercase{\V{lockdown}} \quad // \quad \MakeUppercase{\V{cashout}}}
\chapter{Ending the Game}
\begin{multicols}{2}
\rules{ending}
\end{multicols}

% ================================================================== 08
\renewcommand{\secstatus}{THE LOBBY \quad // \quad THE HOST \quad // \quad SEATS AND BOTS}
\chapter{Hosting, Joining and Leaving}
\begin{multicols}{2}
\rules{hosting}
\end{multicols}

% ================================================================== 09
\renewcommand{\secstatus}{\MakeUppercase{\V{finderscut}} \quad // \quad \MakeUppercase{\V{root}} AND \MakeUppercase{\V{sudo}} \quad // \quad CASH}
\chapter{Where the Money Comes From}
\begin{multicols}{2}
\rules{money}
\end{multicols}

% ================================================================== 10
\renewcommand{\secstatus}{EARLY \quad // \quad MIDDLE \quad // \quad LATE}
\vspace{24pt}\section{Tips}                  % follows on: the money chapter ends early on its page
\begin{multicols}{2}
\rules{tips}
\end{multicols}

% ================================================================== 11
\renewcommand{\secstatus}{THE RULES \quad // \quad THE INTELLIVISION}
\chapter{Frequently Asked Questions}
\begin{multicols}{2}
\rules{faq}

\subsection{On the Intellivision}
\begin{description}[leftmargin=0pt,labelsep=0.5em]
  \item[Which machines?] An Intellivision with a FujiNet Intellivision cartridge, and any TV. An ECS keyboard is
    welcome, not needed. The shot
    clock counts an NTSC console's 60 frames a second; on a PAL console it counts 50, so it stays right.
  \item[Which controller?] The left one. The keypad answers; no overlay is needed, as the prompt line says which
    keys it takes.
  \item[Can I play alone?] Yes. Host a game and start it with \Rule{minplayers} seats: one bot plays you. Any seat
    nobody takes goes to a bot.
  \item[What are the digits on the \V{grid}?] Move hints: each marks a \V{node} that an \V{exploit} in your \V{kit}
    can run on, and the digit is its place in the \V{kit}: yellow if it can run, red if not yet. The options turn
    them on or off.
  \item[Where are the other \Vp{player}?] On the seats view: the disc, left or right, or the lower right button.
    The \V{coin} view shows what everyone holds.
  \item[What do \code{*} and \code{?} before a name mean?] \code{*} is a bot. \code{?} is a \V{player} who is away,
    or on autopilot; a bot plays for them meanwhile.
  \item[Can I turn the sound off?] Yes: in the options, any time. A chord says it's back on.
  \item[Can I play against an Atari, an Apple II, a PC, a CoCo, an NES or an Adam?] Yes. The servers are the same
    for every machine: all of them meet at the same tables.
  \item[Why was my handle refused?] Someone at the table has it, a bot included. If it's your own seat after a
    reset, the server may not have noticed the drop yet: try again in a minute.
  \item[What if everybody leaves?] The game waits for 10 minutes with nobody connected, then it is dropped.
\end{description}
\end{multicols}

% ================================================================== 12
\renewcommand{\secstatus}{EVERY TERM, FROM THE THEME}
\chapter{Glossary}
\rules{glossary}

% ================================================================== REF
\renewcommand{\secstatus}{FIND THE \MakeUppercase{\V{hive}}'S SIZE IN ITS TIER, THEN READ ACROSS}
\clearpage\renewcommand{\secnum}{REF}% after the page break: the running head uses it
\section{\V{coin} Prices and Payouts}
\vspace{6pt}
\begin{center}
\def\textcolor#1#2{\inv{#2}}% the tier's syndicate letters inverse
\newcommand{\hot}[1]{\leavevmode{\color{alarm}#1}}
\machine\fontsize{12}{14}\selectfont
\renewcommand{\arraystretch}{1.3}
\newcolumntype{C}{>{\centering\arraybackslash}p{0.85in}}
\newcolumntype{R}{>{\raggedleft\arraybackslash}p{0.95in}}
\begin{tabular}{CCC@{\hspace{14pt}}R@{\hspace{14pt}}RR}
LOW & MID & HIGH & \leavevmode\color{steel}PRICE & \leavevmode\color{signal}\MakeUppercase{\V{root}} & \leavevmode\color{signal}\MakeUppercase{\V{sudo}} \\
\TierCoins{low} & \TierCoins{mid} & \TierCoins{high} & & & \\[2pt]
\hline
2      & -      & -      & \Y{200}   & \Y{2,000}  & \Y{1,000} \\
3      & 2      & -      & \Y{300}   & \Y{3,000}  & \Y{1,500} \\
4      & 3      & 2      & \Y{400}   & \Y{4,000}  & \Y{2,000} \\
5      & 4      & 3      & \Y{500}   & \Y{5,000}  & \Y{2,500} \\
6-10   & 5      & 4      & \Y{600}   & \Y{6,000}  & \Y{3,000} \\
\hot{11-20}  & 6-10   & 5      & \Y{700}   & \Y{7,000}  & \Y{3,500} \\
21-30  & \hot{11-20}  & 6-10   & \Y{800}   & \Y{8,000}  & \Y{4,000} \\
31-40  & 21-30  & \hot{11-20}  & \Y{900}   & \Y{9,000}  & \Y{4,500} \\
41+    & 31-40  & 21-30  & \Y{1,000} & \Y{10,000} & \Y{5,000} \\
-      & 41+    & 31-40  & \Y{1,100} & \Y{11,000} & \Y{5,500} \\
-      & -      & 41+    & \Y{1,200} & \Y{12,000} & \Y{6,000} \\
\hline
\end{tabular}
\end{center}
\vspace{10pt}
\begin{center}\machine\fontsize{10}{16}\selectfont
\begin{tabular}{@{}ll@{}}
\textcolor{alarm}{RED} & FIRST \MakeUppercase{\V{safe}} SIZE: THE \MakeUppercase{\V{hive}} CAN'T BE \MakeUppercase{\V{absorbed}}\\
\textcolor{signal}{\MakeUppercase{\V{root}}} & 10 x \MakeUppercase{\V{coin}} PRICE\\
\textcolor{signal}{\MakeUppercase{\V{sudo}}} & 5 x \MakeUppercase{\V{coin}} PRICE\\
\Rule{endsize}+ & \MakeUppercase{\V{totalcompromise}}: THE \MakeUppercase{\V{nullroute}} MAY BE CALLED
\end{tabular}\end{center}

\vspace{6pt}
\begin{center}
{\machine\fontsize{9}{11}\selectfont\color{steel} THE LETTERS: EACH SYNDICATE ON THE \MakeUppercase{\V{grid}}, AND ITS TIER}\par\vspace{6pt}
\renewcommand{\arraystretch}{1.3}
\newcommand{\synrow}[1]{\inv{\Syn{#1}{letter}} & \Syn{#1}{name} & {\machine\color{steel}\MakeUppercase{\Syn{#1}{tier}}}}
\begin{tabular}{@{}cl@{\hspace{8pt}}l@{\hspace{36pt}}cl@{\hspace{8pt}}l@{}}
\synrow{0} & \synrow{4} \\
\synrow{1} & \synrow{5} \\
\synrow{2} & \synrow{6} \\
\synrow{3} & & & \\
\end{tabular}
\end{center}

\end{document}
